\chaptertitle{7}{掌握优秀的行业指数基金}


除了策略加权指数基金之外，还有一类基金品种，也是我们重点关注的，那就是优秀行业指数基金。

为何把行业指数基金单独放在后面介绍，而不是和前面的宽基指数基金一起介绍呢？这是因为行业指数的投资风险更高一些，不仅要考虑投资价值，还要考虑不同行业自身的特点和当前所处的发展阶段。

有时候，国家整体的经济发展不错，但有的行业反而一路衰败。

例如过去20年，我们国家经济发展得不错，如果长期投资宽基指数基金，我们可以取得不错的收益。但是如果我们投资了收音机、缝纫机行业，那么随着这些行业的萎缩，我们则会遭受到损失。

行业本身也容易受到一些行业事件和政策变化的影响，例如当年某奶制品污染事件，就对消费行业造成了很大冲击。这都增加了行业指数基金的投资难度。

行业有自己的投资特点，不能一概而论，我们需要针对行业的特点具体分析。所以就投资难度来说，行业指数基金比宽基指数基金要大不少。

那为何还要投资行业指数基金呢？因为行业指数基金能极大地丰富我们的投资品种。有的时候宽基指数基金没有太好的投资机会，但某些行业指数基金反而会出现投资机会，这就能对我们的投资起到很好的补充作用。

不过行业指数基金的投资难度确实大一些，所以刚开始接触指数基金投资的时候，建议投资者从宽基指数基金开始，积累了足够多的经验后再投资行业指数基金。

本章只介绍几种优秀的行业指数及其对应的基金。如果想要了解其他行业指数的知识，可以阅读《指数基金投资指南》一书。

\subsection{哪类行业天生更容易赚钱}
不同行业的赚钱能力是有很大差别的，有的行业天生就更容易赚钱，比如说必需消费和医药行业。

为什么呢？

我们先来想想开公司怎样才能赚到钱。

首先我们的客户得有需求，其次我们的产品要能够满足他们的需求。如果客户的需求稳定，我们才能持续地赚钱，如果客户的需求不稳定，我们就会“饥一顿，饱一顿”。例如很多券商，在牛市的时候门前车水马龙，在熊市的时候门可罗雀。

不过有时候，即使客户需求稳定，我们也不一定能赚到钱。如果生产的产品是同质化的，客户选择我们的产品和选择别人的没有区别，这就是一片“红海”，最终激烈的竞争会导致“价格战”爆发，将利润空间吞噬掉。例如，智能手机行业，很多智能手机生产商的利润很微薄。

那么，客户需求稳定、我们的产品也有别人模仿不了的优点，我们就一定能赚到钱了吗？其实也不然。有的企业需要经常再投入巨资以保证稳定的生产，赚到的钱如果不再投入，企业就经营不下去，这时候我们赚到的钱也不属于自己。例如，某些高污染的重工业、不注重风控的金融业，即使前期赚到许多钱，后期也会还回去。

所以，市场需求比较稳定，企业有“护城河”能保证一定的利润率，再投资需求小，能获得大量的自由现金流，只有拥有这些条件的行业，才能比较容易赚钱。

为什么可口可乐和茅台一度成为价值投资的标签呢？因为它们比较符合上面的条件：需求受经济周期影响小，品牌“护城河”能保证不错的利润率，再投资需求较小，赚到的钱大多都是自由现金流。

如果在二级市场投资这样的企业，想赚钱还要添加一个条件：企业愿意回馈股东，而不是乱花钱。如果企业是赚钱的，但是管理者把赚到的钱随意乱花、铺张浪费，或者进行风险很大的投资，也会导致企业出现问题。

任何一个行业都是国家经济不可缺少的一部分，只要走正确的路，任何一个行业都是能赚钱的，不过赚钱难度明显有差异。

天生更容易赚钱的行业有哪些呢？主要集中在消费和医药行业，例如医药行业、必需消费行业、可选消费行业。

\subsection{必需消费行业：最容易出现大牛股的行业之一}
{\heiti 必需消费行业及其对应的指数基金}

必需消费行业也被称为日常消费行业、主要消费行业。必需消费，主要是维持我们正常生活所需要的各种消费品，例如饮料、酒、农副食品等。

必需消费行业也是需求最稳定的行业，不管经济情况如何，这些日常消费都是我们不可缺少的。正是因为这种很稳定的需求，必需消费行业也是巴菲特最喜欢的行业之一。巴菲特的经典投资案例，很多都是在必需消费行业上。像巴菲特投资的喜诗糖果（See's Candies）、可口可乐、生产番茄酱的亨氏食品（Heinz）、生产剃须刀的吉列（Gillette）等，都是必需消费行业的公司。

{\heiti 必需消费行业指数}

目前，必需消费行业的指数主要是以下四只。

上证消费指数：从上海证券交易所挑选必需消费行业公司。

上证消费80指数：从上海证券交易所挑选80家规模最大的必需消费行业公司。

中证消费指数：从中证800指数，即沪深300指数和中证500指数中挑选必需消费行业公司。

全指消费指数：从所有上市公司中挑选必需消费行业公司，覆盖范围最广。

追踪这些必需消费行业指数的指数基金如表 \ref{tab:7-1} 所示。

{\scriptsize
\setlength{\tabcolsep}{2.75pt}
\renewcommand{\theHtable}{7-1}
\renewcommand{\arraystretch}{1.2}
\renewcommand{\thetable}{7.1}
\begin{longtable}{>{\raggedright\arraybackslash}m{0.28\textwidth}>{\centering\arraybackslash}m{0.13\textwidth}>{\centering\arraybackslash}m{0.12\textwidth}>{\centering\arraybackslash}m{0.15\textwidth}>{\centering\arraybackslash}m{0.12\textwidth}>{\centering\arraybackslash}m{0.12\textwidth}}
\caption{必需消费行业指数基金}\label{tab:7-1}\\[1.2ex]
\hline
\noalign{\vskip 4pt}
\multicolumn{1}{>{\centering\arraybackslash}m{0.28\textwidth}}{\rule[-1.4ex]{0pt}{5.2ex}\bfseries 基金简称} & \multicolumn{1}{>{\centering\arraybackslash}m{0.13\textwidth}}{\bfseries 基金代码} & \multicolumn{1}{>{\centering\arraybackslash}m{0.12\textwidth}}{\shortstack[c]{\bfseries 管理费率\\\bfseries （\%）}} & \multicolumn{1}{>{\centering\arraybackslash}m{0.15\textwidth}}{\shortstack[c]{\bfseries 基金规模\\\bfseries （亿元）}} & \multicolumn{1}{>{\centering\arraybackslash}m{0.12\textwidth}}{\shortstack[c]{\bfseries 成立年限\\\bfseries （年）}} & \multicolumn{1}{>{\centering\arraybackslash}m{0.12\textwidth}}{\bfseries 场内/场外} \\
\hline
\endfirsthead
\multicolumn{6}{r}{续表}\\[1ex]
\hline
\noalign{\vskip 4pt}
\multicolumn{1}{>{\centering\arraybackslash}m{0.28\textwidth}}{\rule[-1.4ex]{0pt}{5.2ex}\bfseries 基金简称} & \multicolumn{1}{>{\centering\arraybackslash}m{0.13\textwidth}}{\bfseries 基金代码} & \multicolumn{1}{>{\centering\arraybackslash}m{0.12\textwidth}}{\shortstack[c]{\bfseries 管理费率\\\bfseries （\%）}} & \multicolumn{1}{>{\centering\arraybackslash}m{0.15\textwidth}}{\shortstack[c]{\bfseries 基金规模\\\bfseries （亿元）}} & \multicolumn{1}{>{\centering\arraybackslash}m{0.12\textwidth}}{\shortstack[c]{\bfseries 成立年限\\\bfseries （年）}} & \multicolumn{1}{>{\centering\arraybackslash}m{0.12\textwidth}}{\bfseries 场内/场外} \\
\hline
\endhead
\hline
\endfoot
\hline
\endlastfoot

华夏消费ETF & 510630 & 0.50 & 1.59 & 6.2 & 场内 \\
汇添富中证主要消费ETF联接 & 000248 & 0.50 & 15.07 & 4.2 & 场内 \\
汇添富中证主要消费ETF & 159928 & 0.50 & 15.25 & 5.8 & 场内 \\
招商上证消费80ETF & 510150 & 0.50 & 1.52 & 8.5 & 场内 \\
嘉实中证主要消费ETF & 512600 & 0.50 & 0.06 & 5.0 & 场内 \\
上投摩根中证消费服务指数 & 370023 & 0.75 & 0.37 & 6.7 & 场外 \\
招商上证消费80ETF联接A & 217017 & 0.50 & 1.25 & 8.5 & 场外 \\
\end{longtable}
}


{\heiti 必需消费行业的主力：食品饮料行业}

必需消费行业的公司90\%都是以食品饮料公司。所以，很多基金公司开发了食品饮料行业指数基金。

食品饮料行业，大概是与我们日常生活联系最紧密的一个行业了。

食品饮料行业中，白酒公司的占比最高，超过50\%。剩下的主要是乳制品、调味品、肉制品等公司。像我们熟悉的伊利、蒙牛、双汇、海天味业等，都属于食品饮料行业的公司。

这些公司跟我们日常生活联系比较紧密，这些消费不太受经济波动的影响。

而且消费习惯一旦养成，是比较难改掉的，比如说北方人逢年过节喜欢喝白酒，小孩子每天喝奶，年轻人经常喝饮料。

另外，成熟的食品饮料公司，本身的利润率都不错，基本都有自己比较知名的品牌，享有一定的品牌溢价。

企业再投入的需求也比较小，也不太需要频繁的升级改造。像智能手机需要投入大量的精力去升级换代，因为消费者买新不买旧，但是像茅台等白酒反而是越老越好，越老越值钱。

巴菲特也是食品饮料行业的忠实爱好者，他很喜欢食品饮料行业“盈利稳定、现金流充沛、不太需要再投入”等特性。

这些特性，也让食品饮料行业成为A股长期收益最高的行业之一。

{\heiti 食品饮料行业指数}

国内的食品饮料行业指数，主要是中证食品饮料行业指数。

这个指数是从2004年年底1000点开始的，到2018年10月，指数点数是11800多点，上涨约11倍。如果加上分红收益，则上涨到14700多点，上涨约14倍。堪称A股收益最高的指数之一。

追踪食品饮料行业的指数基金，主要有：

（1）天弘中证食品饮料指数基金（001631），基金规模2.6亿元，管理费率0.5\%、托管费率0.1\%，平均一年0.6\%的费率。

（2）国泰国证食品饮料指数基金（160222），基金规模17亿元，管理费率1\%、托管费率0.2\%，平均一年1.2\%的费率。

以上两个食品饮料指数基金稍有不同。国证食品饮料指数基金的成份股是50只，中证食品饮料指数基金的成份股是85只，国证食品饮料指数基金更集中于大盘股上。这两只指数基金很相似，估值也差不多。

不过从费率角度，天弘中证食品饮料指数基金的费率低一些，国泰国证食品饮料指数基金是一只分级母基金，费率略高一点。

\subsection{医药行业：经济危机中的避险板块}
{\heiti 医药行业的两大优点：长牛、避险}

每个人都离不开生老病死，生病了一定需要医药。人类这个基本需求不会因为经济不景气、自然灾害等而减少。因此，医药行业天生也是更容易赚钱的行业。

在绝大多数国家，医药行业都是几十年里发展不错的行业。

{\kaishu 以美股为例，从1989年9月11日到2014年7月8日，标普500指数上涨了465\%，而同期标普500医药指数上涨了969\%。医药行业指数的涨幅高出了一倍多，其表现远远超越同期的标普500指数。}

此外，医药行业指数还是经济危机中非常有价值的避险板块。即使整个国家的宏观经济不景气，医药行业指数也很有可能涨得不错。

{\kaishu 以日本股市为例。从1992年到2012年，日本经济陷入“失落的二十年”。日经指数下跌25.6\%，而医药行业指数上涨了92\%，有14年医药行业指数都跑赢了日经指数。}

我国医药行业指数的表现就更突出了。假如从2004年年底开始投资国内的医药行业指数基金，能收获9倍多的收益。

随着国内人口老龄化的严重，医药行业在未来很长一段时间里仍然是持续受益的。以上都说明，投资医药行业指数基金是不错的选择。

不过，相比必需消费行业，医药行业多了一个风险：政策风险。

很少有国家会出台法律去限制可口可乐、啤酒的销售，但几乎所有国家会针对医药行业出台大量的法律法规。

因为大多数时候，国家需要负担国民的一部分医疗支出。随着人口老龄化的严重，医疗支出压力越来越大，因此会反过来限制医药行业的利润空间。政策的变化，会直接影响医药行业的整体业绩。

仍以日本为例，大家都知道，日本老龄化很严重。因此，理论上日本医药行业会大大受益。但是实际上，政府没有足够的钱去支付越来越高的老年人看病的费用。那怎么办呢？

最后，就会演变为日本政府用政策去压迫医药企业，用企业利润来补贴百姓。政策风险会干扰医药行业的定价，从而影响利润率。

不过，优秀行业本身的特性是无法被干扰的。即使加上政策风险，医药行业仍然是数一数二的好行业，收益在所有行业中位于前列。

{\heiti 医药行业指数基金有哪些}

医药行业因为本身比较受欢迎，所以对应的行业指数基金数量也非常多。常被用到的主要有以下几类。

\begin{itemize}
\item 300医药行业指数：挑选了沪深300指数里的医药行业公司，主要是大型医药股。
\item 500医药行业指数：挑选了中证500指数里的医药行业公司，主要是中盘医药股。
\item 中证医药行业指数：包含了300医药指数和500医药指数，基本覆盖了医药行业的大中盘股。
\item 全指医药行业指数：从整个A股中挑选医药行业公司，它覆盖的医药公司是最全的。
\item 中证医药100行业指数：挑选了市值最大的100家公司，等权重分配，每只股票分配比例都只有1\%，定期平衡。
\end{itemize}

有很多基金公司都围绕医药行业指数开发了基金产品，相关指数基金如表 \ref{tab:7-2} 所示。

{\scriptsize
\setlength{\tabcolsep}{2.75pt}
\renewcommand{\theHtable}{7-2}
\renewcommand{\arraystretch}{1.2}
\renewcommand{\thetable}{7.2}
\begin{longtable}{>{\raggedright\arraybackslash}m{0.28\textwidth}>{\centering\arraybackslash}m{0.13\textwidth}>{\centering\arraybackslash}m{0.12\textwidth}>{\centering\arraybackslash}m{0.15\textwidth}>{\centering\arraybackslash}m{0.12\textwidth}>{\centering\arraybackslash}m{0.12\textwidth}}
\caption{医药行业指数基金}\label{tab:7-2}\\[1.2ex]
\hline
\noalign{\vskip 4pt}
\multicolumn{1}{>{\centering\arraybackslash}m{0.28\textwidth}}{\rule[-1.4ex]{0pt}{5.2ex}\bfseries 基金简称} & \multicolumn{1}{>{\centering\arraybackslash}m{0.13\textwidth}}{\bfseries 基金代码} & \multicolumn{1}{>{\centering\arraybackslash}m{0.12\textwidth}}{\shortstack[c]{\bfseries 管理费率\\\bfseries （\%）}} & \multicolumn{1}{>{\centering\arraybackslash}m{0.15\textwidth}}{\shortstack[c]{\bfseries 基金规模\\\bfseries （亿元）}} & \multicolumn{1}{>{\centering\arraybackslash}m{0.12\textwidth}}{\shortstack[c]{\bfseries 成立年限\\\bfseries （年）}} & \multicolumn{1}{>{\centering\arraybackslash}m{0.12\textwidth}}{\bfseries 场内/场外} \\
\hline
\endfirsthead
\multicolumn{6}{r}{续表}\\[1ex]
\hline
\noalign{\vskip 4pt}
\multicolumn{1}{>{\centering\arraybackslash}m{0.28\textwidth}}{\rule[-1.4ex]{0pt}{5.2ex}\bfseries 基金简称} & \multicolumn{1}{>{\centering\arraybackslash}m{0.13\textwidth}}{\bfseries 基金代码} & \multicolumn{1}{>{\centering\arraybackslash}m{0.12\textwidth}}{\shortstack[c]{\bfseries 管理费率\\\bfseries （\%）}} & \multicolumn{1}{>{\centering\arraybackslash}m{0.15\textwidth}}{\shortstack[c]{\bfseries 基金规模\\\bfseries （亿元）}} & \multicolumn{1}{>{\centering\arraybackslash}m{0.12\textwidth}}{\shortstack[c]{\bfseries 成立年限\\\bfseries （年）}} & \multicolumn{1}{>{\centering\arraybackslash}m{0.12\textwidth}}{\bfseries 场内/场外} \\
\hline
\endhead
\hline
\endfoot
\hline
\endlastfoot

易方达沪深300医药卫生ETF联接 & 001344 & 0.50 & 2.30 & 1.5 & 场外 \\
易方达沪深300医药ETF & 512010 & 0.50 & 4.74 & 5.7 & 场内 \\
南方中证500医药卫生ETF & 512300 & 0.50 & 0.40 & 4.6 & 场内 \\
广发医药卫生联接A & 001180 & 0.50 & 9.64 & 4.1 & 场外 \\
广发医药卫生联接C & 002978 & 0.50 & 0.56 & 2.9 & 场外 \\
广发中证全指医药卫生ETF & 159938 & 0.50 & 16.95 & 4.5 & 场内 \\
华夏医药ETF & 510660 & 0.50 & 0.71 & 6.2 & 场内 \\
天弘中证医药100A & 001550 & 0.50 & 2.04 & 3.9 & 场外 \\
天弘中证医药100C & 001551 & 0.50 & 1.64 & 3.9 & 场外 \\
汇添富中证医药卫生ETF & 159929 & 0.50 & 1.23 & 5.8 & 场内 \\
嘉实中证医药卫生ETF & 512610 & 0.50 & 0.11 & 5.0 & 场内 \\
鹏华中证医药卫生（LOF） & 160635 & 0.75 & 0.48 & 3.8 & 场内/场外 \\
国联安中证医药100A & 000059 & 0.80 & 14.76 & 5.8 & 场外 \\
国联安中证医药100C & 006569 & 0.80 & 0.04 & 0.6 & 场外 \\
国泰国证医药卫生行业指数分级 & 160219 & 1.00 & 19.02 & 5.7 & 场外 \\
\end{longtable}
}


不过，并不是所有的医药行业指数基金都适合投资。有的基金规模太小，就不太合适。下面螺丝钉具体分析一下每个品类指数基金的特点。

\begin{itemize}
\item 300医药行业指数基金：优点是，规模大，场内、场外都能买，有利于头部医药股的发挥。缺点是，如果头部医药股出问题了，那受到的影响也比较大。属于高风险、高收益的医药行业指数基金。
\item 500医药行业指数基金：规模较小，不到1亿元，我们通常不考虑这样的品种。
\item 中证医药行业指数基金：跟500医药行业指数基金一样，大部分规模也不大，只有汇添富中证医药卫生ETF规模符合我们的要求。
\item 全指医药行业指数基金：广发基金的全指医药基金规模比较合适。
\item 中证医药100行业指数基金：国联安和天弘医药100比较合适，其中天弘医药100的费率比国联安的要低一些，如果看好中证医药100行业指数，可以考虑天弘医药100指数基金。
\end{itemize}

那具体投资的时候，该选择哪一只医药行业指数基金呢？

其实从收益率的角度，这几只医药行业指数基金都还不错。不过，很容易出现“黑天鹅”风险，也容易受到政策影响。假如指数基金过于重仓某些医药股，其实还是有风险的。

如果能做到足够分散，例如像中证医药100行业指数基金这样，每只成份股都只有1\%左右的比例，那么，即使指数中某些股票出了问题，但是对整体的影响并不大。

如果愿意承担一些风险，可以选择全指医药行业指数基金，或者300医药行业指数基金。

\subsection{可选消费行业：升级换代快，有周期性}
{\heiti 什么是可选消费行业}

必需消费指的是我们在日常生活中经常需要的食品、饮料等消费品，一般单价较低，消费频率较高，也是刚需。而可选消费则是和必需消费对应的。

什么是可选消费呢？从名字上也可以大致区分出来，可选消费指的是有钱的时候才会消费，这种消费可以提升我们的生活质量，但不是刚需。没有钱的时候，我们会推迟可选消费方面的支出，像高档手机、汽车、大家电等，都属于可选消费的范畴。

{\heiti 可选消费的特点}

可选消费有如下特点。

（1）需求比必需消费弱，有一定周期性。

我们可以一星期不用空调、不开汽车，但是估计没有人能一星期不吃饭、不喝水。可选消费的需求不像必需消费那样稳定，如果经济形势不好，消费者会倾向于将可选消费推后，或者降级选用更便宜的产品。所以，很多时候可选消费会表现出更强的周期性。像家用电器、汽车行业等，都存在很强的业绩周期波动。

（2）受益于人口红利，特别是人均消费金额的提升。

可选消费非常受益于人口红利，特别是人均消费金额的提升。人口的增长、人均消费金额达到一定程度，都会促进对应的可选消费企业的兴起。

例如影视行业，这几年影视行业的爆发式发展大家有目共睹，但影视行业在之前很长一段时期都处于低速发展期。几年前，全国影视行业的利润加起来还不如一家排名第10的银行的利润，但现在票房超过10亿元的大片比比皆是。虽然这其中有网络发展、电影产业成熟等因素的推动，但最根本的原因还是人均消费金额的提升。

由于目前国内人口老龄化在加剧，新增人口数在短期内来看并不乐观，所以可选消费未来的增长点主要在于人均消费金额的提升。我国人均每日消费金额的中位数在40～50元，这意味着还有7亿多居民的人均每日消费金额不足50元，未来还是有很大提升空间的。

（3）具有升级换代的特性。

可选消费有升级换代的特性。每个时代都有其具有代表性的可选消费产品和企业，这些企业在自己所引领的时代里能呼风唤雨，股价和业绩表现都非常惊人，但如果跟不上消费升级换代，企业就会衰落。

这一点跟必需消费形成了鲜明对比。我们对饮料、油盐酱醋等必需消费产品的升级换代并没有过多要求，今天吃的盐跟20年前吃的盐相比，差别并不大。但每隔一段时间，可选消费的产品就会升级换代。

这个特点会导致一个问题，就是如果某只可选消费行业的个股，在升级换代中被替换掉了，那么对投资者来说就会有很大的风险。所以对普通投资者来说，投资可选消费行业的指数是个不错的选择，指数囊括了一系列的相关股票，可以避免投资个股的风险。

目前可选消费行业指数基金比较少，如表 \ref{tab:7-3} 所示。

{\scriptsize
\setlength{\tabcolsep}{2.75pt}
\renewcommand{\theHtable}{7-3}
\renewcommand{\arraystretch}{1.2}
\renewcommand{\thetable}{7.3}
\begin{longtable}{>{\raggedright\arraybackslash}m{0.28\textwidth}>{\centering\arraybackslash}m{0.13\textwidth}>{\centering\arraybackslash}m{0.12\textwidth}>{\centering\arraybackslash}m{0.15\textwidth}>{\centering\arraybackslash}m{0.12\textwidth}>{\centering\arraybackslash}m{0.12\textwidth}}
\caption{可选消费行业指数基金}\label{tab:7-3}\\[1.2ex]
\hline
\noalign{\vskip 4pt}
\multicolumn{1}{>{\centering\arraybackslash}m{0.28\textwidth}}{\rule[-1.4ex]{0pt}{5.2ex}\bfseries 基金简称} & \multicolumn{1}{>{\centering\arraybackslash}m{0.13\textwidth}}{\bfseries 基金代码} & \multicolumn{1}{>{\centering\arraybackslash}m{0.12\textwidth}}{\shortstack[c]{\bfseries 管理费率\\\bfseries （\%）}} & \multicolumn{1}{>{\centering\arraybackslash}m{0.15\textwidth}}{\shortstack[c]{\bfseries 基金规模\\\bfseries （亿元）}} & \multicolumn{1}{>{\centering\arraybackslash}m{0.12\textwidth}}{\shortstack[c]{\bfseries 成立年限\\\bfseries （年）}} & \multicolumn{1}{>{\centering\arraybackslash}m{0.12\textwidth}}{\bfseries 场内/场外} \\
\hline
\endfirsthead
\multicolumn{6}{r}{续表}\\[1ex]
\hline
\noalign{\vskip 4pt}
\multicolumn{1}{>{\centering\arraybackslash}m{0.28\textwidth}}{\rule[-1.4ex]{0pt}{5.2ex}\bfseries 基金简称} & \multicolumn{1}{>{\centering\arraybackslash}m{0.13\textwidth}}{\bfseries 基金代码} & \multicolumn{1}{>{\centering\arraybackslash}m{0.12\textwidth}}{\shortstack[c]{\bfseries 管理费率\\\bfseries （\%）}} & \multicolumn{1}{>{\centering\arraybackslash}m{0.15\textwidth}}{\shortstack[c]{\bfseries 基金规模\\\bfseries （亿元）}} & \multicolumn{1}{>{\centering\arraybackslash}m{0.12\textwidth}}{\shortstack[c]{\bfseries 成立年限\\\bfseries （年）}} & \multicolumn{1}{>{\centering\arraybackslash}m{0.12\textwidth}}{\bfseries 场内/场外} \\
\hline
\endhead
\hline
\endfoot
\hline
\endlastfoot

广发中证全指可选消费ETF联接A & 001133 & 0.50 & 2.33 & 4.1 & 场外 \\
广发中证全指可选消费ETF & 159936 & 0.50 & 2.25 & 5.0 & 场内 \\
\end{longtable}
}


\subsection{行业指数基金该如何选择}

从长期收益的角度，消费、医药行业指数基金是高于所有行业指数基金平均值的，自然也是高于沪深300和中证500等普通指数基金的，跟部分策略加权指数基金的长期收益是一个级别。

那这么多的优秀行业指数基金，我们该如何选择呢？

行业指数本身收益好，但是相比宽基指数来说，波动也更大，也更容易受到政策、行业意外事件的影响。例如消费行业就出现过的一些负面事件，医药行业也出现过针对行业整体的政策限制。

投资这些行业指数基金时，需要承担这些额外的风险，所以投资难度比一般的宽基指数基金要大一些。

通常我们定投时，可以以优秀的宽基策略加权指数基金为主，以优秀行业指数基金为辅。例如以红利、基本面、价值、低波动四类策略加权指数基金为主，以消费、医药、中概互联行业指数基金为辅。

宽基指数基金比较稳健，适合作为投资主力。优秀行业指数基金长期收益出色，可以提高我们的定投收益效果。这样的思路，更适合普通投资者。

这几章，我们都是从“买什么”的角度，来优化我们的定投。

我们要挑选规模合适、费率低廉、追踪误差小的指数基金。

我们要挑选长期收益出色的策略加权指数基金和优秀行业指数基金。

定投还有两个很重要的环节，就是“怎么买”和“怎么卖”的问题。接下来，我们就来了解一下这两个环节该怎么优化和提高收益。

{\heiti 投资者笔记}

\begin{itemize}
\item {\kaishu 天生更容易赚钱的行业主要集中在消费和医药行业。}
\item {\kaishu 宽基指数基金比较稳健，适合作为投资主力。优秀行业指数基金长期收益出色，可以提高我们的定投收益效果。}
\end{itemize}

